\documentclass[11pt]{article}
\usepackage[margin=1in]{geometry}
\usepackage{amsmath,amssymb,amsthm,mathtools}
\usepackage[colorlinks=true,linkcolor=blue,citecolor=blue,urlcolor=blue]{hyperref}

\newcommand{\dd}{\,\mathrm{d}}
\newcommand{\Ree}{\operatorname{Re}}
\newcommand{\Tr}{\operatorname{Tr}}
\newcommand{\FP}{\operatorname{FP}}

\title{A clean route to zeta regularization:\\
\large heat kernel, Mellin transform, and meromorphic continuation}
\author{P. Shubham Parashar}
\date{Cleaned working note}

\begin{document}
\maketitle

\begin{abstract}
This note separates three distinct operations that are often conflated in informal discussions of zeta regularization: ordinary summation, regulator-dependent finite parts, and analytic continuation. The basic example is the exponential cutoff of $\sum_{n\ge 1} n$, followed by the Mellin-transform construction of the Riemann zeta function and its spectral generalization for Laplace-type operators on compact manifolds.
\end{abstract}

\section{Three notions that must be kept separate}
There are three different objects in play.
\begin{itemize}
  \item An \emph{ordinary sum} $\sum_n a_n$ means the limit of partial sums, when that limit exists.
  \item A \emph{regularization} replaces a divergent expression by a convergent family $F(\varepsilon)$, studies the asymptotics as $\varepsilon\to 0^+$, and may extract a finite part.
  \item An \emph{analytic continuation} extends a holomorphic function from its initial domain to a larger domain, typically as a meromorphic function.
\end{itemize}
The slogan $1+2+3+\cdots=-1/12$ is therefore misleading unless it is explicitly interpreted as a statement about a particular regularized finite part or about the analytically continued value of the Riemann zeta function at $s=-1$.

\section{A convergent cutoff calculation}
For $\varepsilon>0$, define
\begin{equation}
  S(\varepsilon)=\sum_{n=1}^{\infty} n e^{-\varepsilon n}.
\end{equation}
This is an ordinary convergent sum for each positive $\varepsilon$. With $q=e^{-\varepsilon}$,
\begin{equation}
  \sum_{n\ge1} n q^n = \frac{q}{(1-q)^2},
\end{equation}
so
\begin{equation}
  S(\varepsilon)=\frac{e^{-\varepsilon}}{(1-e^{-\varepsilon})^2}.
\end{equation}
The Bernoulli expansion gives
\begin{equation}
  \frac{\varepsilon}{1-e^{-\varepsilon}}
  =1+\frac{\varepsilon}{2}+\frac{\varepsilon^2}{12}-\frac{\varepsilon^4}{720}+O(\varepsilon^6),
\end{equation}
and hence
\begin{equation}
  S(\varepsilon)=\frac{1}{\varepsilon^2}-\frac{1}{12}+O(\varepsilon^2)
  \qquad (\varepsilon\to0^+).
\end{equation}
The divergent part is $\varepsilon^{-2}$ and the cutoff-independent constant term is $-1/12$.

\paragraph{Finite part in this scheme.}
For this exponential cutoff one may define
\begin{equation}
  \sum_{n\ge1}^{(\mathrm{FP})} n
  :=\FP_{\varepsilon\to0^+}\sum_{n=1}^{\infty}n e^{-\varepsilon n}
  =-\frac{1}{12},
\end{equation}
where $\FP$ means the constant term in the small-$\varepsilon$ asymptotic expansion after divergent powers have been removed. This is not an ordinary sum.

\section{Mellin transforms as the systematic mechanism}
For a function $F(t)$ on $(0,\infty)$, its Mellin transform is
\begin{equation}
  \mathcal{M}[F](s)=\int_0^{\infty} t^{s-1}F(t)\dd t,
\end{equation}
where the integral converges. The important structural fact is that small-$t$ asymptotics of $F(t)$ control the pole structure of its Mellin transform.

Suppose that as $t\to0^+$,
\begin{equation}
  F(t)\sim \sum_j c_j t^{\alpha_j}.
\end{equation}
Then the terms $c_j t^{\alpha_j}$ generate possible simple poles at $s=-\alpha_j$ in the Mellin transform. The standard continuation procedure is to subtract finitely many singular small-$t$ terms, integrate the subtracted expression in a larger strip, and then add back the subtracted terms as explicit meromorphic functions of $s$.

\section{Riemann zeta from a heat-kernel type integral}
For $\Ree(s)>1$,
\begin{equation}
  \zeta(s)=\sum_{n=1}^{\infty} n^{-s}.
\end{equation}
Using
\begin{equation}
  \frac{1}{e^t-1}=\sum_{n=1}^{\infty} e^{-nt}
\end{equation}
and
\begin{equation}
  \int_0^{\infty} t^{s-1}e^{-nt}\dd t=\Gamma(s)n^{-s},
\end{equation}
one obtains, again for $\Ree(s)>1$,
\begin{equation}
  \zeta(s)=\frac{1}{\Gamma(s)}\int_0^{\infty}\frac{t^{s-1}}{e^t-1}\dd t.
\end{equation}
Near $t=0$,
\begin{equation}
  \frac{1}{e^t-1}=\frac{1}{t}-\frac{1}{2}+\frac{t}{12}-\frac{t^3}{720}+\cdots.
\end{equation}
Subtracting finitely many terms of this expansion and adding them back explicitly produces the meromorphic continuation of $\zeta(s)$. In that analytically continued sense,
\begin{equation}
  \zeta(-1)=-\frac{1}{12}.
\end{equation}
This value agrees with the finite part of the exponential cutoff above, but the two statements should still be distinguished conceptually.

\section{Spectral zeta functions and heat traces}
Let $M$ be a compact Riemannian manifold and let $\Delta\ge0$ be the Laplace--Beltrami operator. Its spectrum is discrete,
\begin{equation}
  0=\lambda_0\le\lambda_1\le\lambda_2\le\cdots,
  \qquad \lambda_k\to\infty.
\end{equation}
The heat trace is the convergent quantity
\begin{equation}
  \Theta(t)=\Tr(e^{-t\Delta})=\sum_{k\ge0}e^{-t\lambda_k},
  \qquad t>0.
\end{equation}
As $t\to0^+$ it has an asymptotic expansion
\begin{equation}
  \Theta(t)\sim (4\pi t)^{-\dim M/2}\sum_{j\ge0}a_j t^{j/2}.
\end{equation}
The spectral zeta function is initially defined for $\Ree(s)$ large by
\begin{equation}
  \zeta_{\Delta}(s)=\sum_{k\ge1}\lambda_k^{-s}.
\end{equation}
Removing the zero mode from the heat trace gives the Mellin representation
\begin{equation}
  \zeta_{\Delta}(s)=\frac{1}{\Gamma(s)}\int_0^{\infty}t^{s-1}\bigl(\Theta(t)-1\bigr)\dd t,
\end{equation}
where this integral converges. The heat-kernel expansion at small $t$ then generates the meromorphic continuation of $\zeta_{\Delta}(s)$ to the complex plane. The coefficients $a_j$ determine the possible pole structure.

\section{A good sentence and a bad sentence}
A careful statement is:
\begin{quote}
In the exponential cutoff scheme, the finite part of $\sum_{n\ge1}n e^{-\varepsilon n}$ as $\varepsilon\to0^+$ is $-1/12$; equivalently, the analytically continued Riemann zeta function satisfies $\zeta(-1)=-1/12$.
\end{quote}
A bad statement is:
\begin{quote}
$1+2+3+\cdots=-1/12$.
\end{quote}
The second phrase erases the distinction between ordinary summation, regularization, and analytic continuation.

\end{document}
